% --- IMPORTS ---
\documentclass[leqno]{article}
\usepackage{graphicx}
\usepackage{dsfont}
\usepackage{mathtools}
\usepackage{amssymb}
\usepackage{amsmath}
\usepackage{amsthm}
\usepackage{float}
\usepackage{ragged2e}
\usepackage{tensor}
\usepackage{fancyhdr}
\usepackage{empheq}
\usepackage[most]{tcolorbox}
\usepackage{enumitem}
\usepackage{dirtytalk}
\usepackage{marginnote}
\usepackage{microtype}
\usepackage[
  a4paper,
  left=0.8in,
  right=1in,
  top=1in,
  bottom=0.5in,
  headheight=14pt,
  headsep=0.4in
]{geometry}
\setlength{\parindent}{0cm}
% --- TikZ Packages ---
\usepackage{pgfplots}
\pgfplotsset{compat=1.18}
\usepgfplotslibrary{fillbetween, polar}
\usetikzlibrary{calc, positioning}
\usetikzlibrary{angles, quotes}
\usetikzlibrary{arrows.meta}
\usetikzlibrary{decorations.pathreplacing, decorations.markings}
\usetikzlibrary{patterns}
\usetikzlibrary{intersections}
% --- URL Packages ---
\usepackage[colorlinks=true,linkcolor=red,citecolor=magenta,urlcolor=black]{hyperref}%
\urlstyle{same}
% --- END OF IMPORTS ---

% --- CUSTOM COMMANDS ---
\RenewDocumentCommand{\marks}{o m}{%
  \IfValueTF{#1}%
    {\marginnote{\raggedright\textbf{[#2]}}[#1]}%
    {\marginnote{\raggedright\textbf{[#2]}}}%
}
\newcommand{\vspacerule}[2]{%
  \par\vspace*{#1}%
  \noindent\hrulefill\par
  \vspace*{#2}%
}
\definecolor{scarlet}{RGB}{205, 40, 75}
\newcommand{\questionitem}{%
  \item\phantomsection
  \addcontentsline{toc}{section}{Question \number\value{enumi}}%
}
\renewcommand{\contentsname}{Questions}
% --- END OF CUSTOM COMMANDS ---

% --- HEADER CONFIG ---
\pagestyle{fancy}
\fancyhead{}
\fancyfoot{}
\renewcommand{\headrulewidth}{0.9pt}
\lhead{Integration with Inverse Trig}
\chead{}
\rhead{\href{https://danieljsmith.org}{danieljsmith.org}}
% --- END OF HEADER CONFIG ---

\begin{document}
\vspace*{20pt}
\tableofcontents
\newpage
\begin{enumerate}[
  label=\textbf{\arabic*.},
  leftmargin=*,
  topsep=0pt,
  partopsep=0pt,
  parsep=0pt
]

% ---- Question 1 ----
\questionitem

Find the exact value of the improper integral
\[
\int_0^4 \frac{1}{\sqrt{x(4-x)}}\,\mathrm{d}x
\]
\marks[-30pt]{5}

\newpage

% ---- Question 2 ----
\questionitem

Show that
\[
\int_{0}^{3} \frac{1}{x^2 - 2x + 5}\,\mathrm{d}x = \frac{1}{2}\arctan (3)
\]
\marks[-20pt]{5}

\newpage

% ---- Question 3 ----
\questionitem

\begin{enumerate}[label=(\roman*)]
  \item Prove, from definitions involving exponentials, that
  \[
  1-\tanh^2u=\operatorname{sech}^2u
  \]
  \marks[-20pt]{3}
  \item Prove that, for $|y|<1$,
  \[
  \operatorname{artanh}y=\frac{1}{2}\ln\left(\frac{1+y}{1-y}\right)
  \]
  \marks[-20pt]{4}
  \item Use the substitution $x=4\tanh u$ to show that, for $|x|<4$,
  \[
  \int \frac{1}{16-x^2}\,\mathrm{d}x=\frac{1}{4}\operatorname{artanh}\frac{x}{4}+c
  \]
  Hence deduce that
  \[
  \int \frac{1}{16-x^2}\,\mathrm{d}x=\frac{1}{8}\ln\left(\frac{4+x}{4-x}\right)+c
  \]
  where $c$ is an arbitrary constant. \marks{6} \\
  \item By first expressing $-t^2+6t+7$ in completed square form, show that
  \[
  \int_{1}^{5} \frac{1}{-t^2+6t+7}\,\mathrm{d}t=\frac{1}{4}\ln 3
  \]
  \marks[-20pt]{5}
\end{enumerate}
\newpage
% ---- Question 4 ----
\questionitem

\begin{enumerate}[label=(\roman*), start=1]
  \item For $x>5$, find
  \[
  \int \frac{1}{\sqrt{x^2-6x+5}}\,\mathrm{d}x
  \] \\
  expressing your answer first in terms of $\operatorname{arcosh}$ and hence in logarithmic form. \marks{3} \\
  \item Use integration by substitution with $x-3=2\cosh u$ to show that, for $x>5$,
  \[
  \int \sqrt{x^2-6x+5}\,\mathrm{d}x=2\left(\frac{x-3}{2}\sqrt{\frac{(x-3)^2}{4}-1}-\ln\left|\frac{x-3}{2}+\sqrt{\frac{(x-3)^2}{4}-1}\right|\right)+c
  \]
  where $c$ is an arbitrary constant. \marks{7} \\
\end{enumerate}

\newpage

% ---- Question 5 ----
\questionitem

\[
g(x)=\frac{2x^2-3x+31}{x^3-3x^2+9x+13}
\]
\\
\begin{enumerate}[label=\textbf{(\alph*)}]
  \item Express $g(x)$ in the form \[\frac{P}{x+1}+\frac{Q}{x^2-4x+13}\] where $P$ and $Q$ are constants to be found. \marks{3} \\
  \item Find $\int g(x)\,\mathrm{d}x$, giving your answer in the form \[A\ln|x+1|+B\arctan\left(\frac{x-2}{3}\right)+c\] where $A$ and $B$ are constants to be found. \marks{4} \\
  \item Hence show that $\int_0^\infty g(x)\,\mathrm{d}x$ diverges. \marks{2} \\
\end{enumerate}


\newpage

% ---- Question 6 ----
\questionitem

Show that
\[
\int_{1}^{4} \frac{1}{\sqrt{10 + 8x - 2x^2}}\,\mathrm{d}x = \frac{1}{\sqrt{2}}\left(\arcsin \left(\frac{2}{3}\right) + \arcsin \left(\frac{1}{3}\right)\right)
\]
\marks[-20pt]{6}

\newpage

% ---- Question 7 ----
\questionitem

\begin{enumerate}[label=\textbf{(\alph*)}]
  \item Explain why
  \[
  \int_{0}^{\infty} \frac{1}{x^2 - 6x + 18} \, \mathrm{d}x
  \]
  is an improper integral. \marks{1} \\
  \item Show that
  \[
  \int_{0}^{\infty} \frac{1}{x^2 - 6x + 18} \, \mathrm{d}x = k\pi
  \]
  where $k$ is a constant to be determined. \marks{4} \\
\end{enumerate}

\newpage

% ---- Question 8 ----
\questionitem

Evaluate
\[
\int_{0}^{4} \frac{1}{x^2-4x+8}\, \mathrm{d}x
\]
giving your answer as an exact multiple of $\pi$. \marks{8}

\newpage

% ---- Question 9 ----
\questionitem

\[f(x)=\frac{6x+1}{4x^2-8x+13}\] \\

\begin{enumerate}[label=\textbf{(\alph*)}]
  \item Find $\int f(x)\,\mathrm{d}x$, giving your answer in the form \[A\ln(4x^2-8x+13)+B\arctan\left(\frac{2x-2}{3}\right)+c\] where $c$ is an arbitrary constant and $A$ and $B$ are constants to be found. \marks{4} \\
  \item Hence show that the mean value of $f(x)$ over the interval $\left[1,\frac{5}{2}\right]$ is \[\frac{1}{2}\ln 2+\frac{7\pi}{36}\] \marks[-20pt]{3} 
  \item Hence write down the mean value of $3-2f(x)$ over the interval $\left[1,\frac{5}{2}\right]$. \marks{1} \\
\end{enumerate}

\newpage

% ---- Question 10 ----
\questionitem

\[
g(x)=\frac{4x-1}{x^2-4x+8}
\] \\
Show that
\[
\int g(x)\,\mathrm{d}x=A\arctan\left(\frac{x-2}{2}\right)+B\ln(x^2-4x+8)+c
\]
where $c$ is an arbitrary constant and $A$ and $B$ are constants to be found. \marks{5} \\
\newpage

% ---- Question 11 ----
\questionitem

You are given that $y=\operatorname{arcosh}x$, where $x\ge 1$. \\

\begin{enumerate}[label=(\roman*), start=1]
  \item Show that
  \[
  \frac{\mathrm{d}y}{\mathrm{d}x}=\frac{1}{\sqrt{x^2-1}}
  \]
  \marks[-20pt]{4}
  \item Show, by integrating the result in part (i), that
  \[
  y=\ln\left(x+\sqrt{x^2-1}\right)
  \]
  \marks[-20pt]{4}
  \item Hence show that
  \[
  \int_{15/8}^{5/2} \frac{1}{\sqrt{4x^2-9}}\,\mathrm{d}x=\frac{1}{2}\left(\operatorname{arcosh}\frac{5}{3}-\operatorname{arcosh}\frac{5}{4}\right)
  \]
  Express this answer in logarithmic form. \marks{4} \\
\end{enumerate}

\newpage


% ---- Question 12 ----
\questionitem

Determine each of the following integrals. \\
\begin{enumerate}[label=(\roman*)]
  \item
  \[
  \int \frac{1}{9x^2+12x+20}\,\mathrm{d}x
  \]
  \marks[-20pt]{4}
  \item
  \[
  \int \frac{1}{\sqrt{45+12x-9x^2}}\,\mathrm{d}x
  \]
  \marks[-20pt]{4}
\end{enumerate}

\newpage

% ---- Question 13 ----
\questionitem

Without using a calculator, find \\
\begin{enumerate}[label=\textbf{(\alph*)}]
  \item
  \[
  \int_{-1/5}^{1} \frac{1}{25x^2-20x+13}\,\mathrm{d}x
  \] \\
  giving your answer as a multiple of $\pi$ \marks{5} \\
  \item
  \[
  \int_{-1/5}^{1} \frac{1}{\sqrt{25x^2-20x+13}}\,\mathrm{d}x
  \] \\
  giving your answer in the form $p\ln(q+r\sqrt{2})$, where $p$, $q$ and $r$ are rational numbers to be found. \marks{7} \\
\end{enumerate}

\newpage

% ---- Question 14 ----
\questionitem

\begin{enumerate}[label=(\roman*)]
  \item Given that $\tanh y=x$, where $|x|<1$, show that
  \[
  y=\frac{1}{2}\ln\left(\frac{1+x}{1-x}\right)
  \]
  and hence that
  \[
  \operatorname{artanh}x=\frac{1}{2}\ln\left(\frac{1+x}{1-x}\right)
  \]
  \marks[-20pt]{6}
  \item Use your result in part (i), or otherwise, to find
  \[
  \int_{1/5}^{2/5} \frac{1}{4-9x^2}\,\mathrm{d}x
  \]
  expressing your answer in an exact logarithmic form. \marks{6} \\
\end{enumerate}

\newpage

% ---- Question 15 ----
\questionitem

Show that
\[
\int \frac{x^2+14x+21}{(x-1)(x^2+4x+13)}\,\mathrm{d}x=A\ln|x-1|+B\ln(x^2+4x+13)+D\arctan\left(\frac{x+2}{3}\right)+c
\]
where $A$, $B$ and $D$ are constants to be found. \marks{7} \\

\newpage

% ---- Question 16 ----
\questionitem

\begin{figure}[H]
    \centering
\begin{tikzpicture}[scale=1.3]
\def\xmin{-2.2}
\def\xmax{5.2}
\def\ymin{-0.3}
\def\ymax{2.8}
\def\num{9}
\def\peak{1}
\def\den{4}
\def\xL{-1}
\def\xR{4}
\pgfmathsetmacro{\yL}{\num/((\xL-\peak)*(\xL-\peak)+\den)}
\pgfmathsetmacro{\yR}{\num/((\xR-\peak)*(\xR-\peak)+\den)}
\coordinate (O) at (0,0);
\coordinate (L) at (\xL,0);
\coordinate (R) at (\xR,0);
\coordinate (LT) at (\xL,\yL);
\coordinate (RT) at (\xR,\yR);
\draw[thick,-{Stealth[length=3mm]}] (\xmin,0) -- (\xmax,0) node[right] {$x$};
\draw[thick,-{Stealth[length=3mm]}] (0,\ymin) -- (0,\ymax) node[above] {$y$};
\draw[thick] plot[domain=\xL:\xR,samples=200] (\x,{\num/((\x-\peak)*(\x-\peak)+\den)});
\draw[thick] (L) -- (LT);
\draw[thick] (R) -- (RT);
\node[below] at (L) {$-1$};
\node[below] at (R) {$4$};
\node[below left] at (O) {$O$};
\end{tikzpicture}
\end{figure}


The diagram shows the cross-section of a solid wooden beam. \\
Each unit on the axes represents $6$ cm. The curved upper boundary of the beam is modelled by the equation
\[
y=\frac{9}{(x-1)^2+4}
\] \\
The cross-section is bounded by the curve, the $x$-axis and the lines $x=-1$ and $x=4$. \\
Assuming that the beam is a prism of length $240$ cm, find, correct to $3$ significant figures, the volume of the beam. \marks{6} \\

\newpage

% ---- Question 17 ----
\questionitem

Find
\[
\int_1^{e} \frac{1}{x\sqrt{4-(\ln x)^2}} \, \mathrm{d}x
\]
expressing your answer in terms of $\pi$. \marks{8}

\end{enumerate}
\end{document}
